% 张三慧《大学物理学》第4版第二卷，第15章《静电场中的电介质》
% 原教材扫描PDF物理页99--102；原印刷页82--85。
% 区别第15章“思考题”和“习题”，后者原题号15.1--15.30。
% 候选转写稿，图15.20--15.29的几何和所有答案仍待专家独立审核。
\section*{第15章\quad 静电场中的电介质}
\subsection*{章末习题（15.1--15.30）}

\paragraph{习题 15.1}
在 \(\mathrm{HCl}\) 分子中，氯核和质子（氢核）的距离为
\(0.128\,\mathrm{nm}\)。假设氢原子的电子完全转移到氯原子上，
并与其他电子构成球对称的负电荷分布，其中心就在氯核上。
此模型的电矩多大？实测的 HCl 分子的电矩为
\(3.4\times10^{-30}\,\mathrm{C\,m}\)，
HCl 分子中的负电分布的「重心」应在何处？
（氯核的电量为 \(17e\)。）

\paragraph{习题 15.2}
两个同心的薄金属球壳，内、外球壳半径分别为
\(R_1=0.02\,\mathrm m\) 和 \(R_2=0.06\,\mathrm m\)。
球壳间充满两层均匀电介质，其相对介电常量分别为
\(\varepsilon_{r1}=6\) 和 \(\varepsilon_{r2}=3\)，
两层介质的分界面半径 \(R=0.04\,\mathrm m\)。
设内球壳带电量 \(Q=-6\times10^{-8}\,\mathrm C\)，求：
（1）\(\boldsymbol D,\boldsymbol E\) 的分布并画 \(D-r,E-r\) 曲线；
（2）两球壳之间的电势差；
（3）贴近内金属壳的电介质表面上的面束缚电荷密度。

\paragraph{习题 15.3}
两共轴的导体圆筒的内、外筒半径分别为
\(R_1,R_2\)，\(R_2<2R_1\)。
其间有两层均匀介质，分界面半径为 \(r_0\)。
内层介质相对介电常量为 \(\varepsilon_{r1}\)，
外层为 \(\varepsilon_{r2}\)，且
\(\varepsilon_{r2}=\varepsilon_{r1}/2\)。
两层介质的击穿场强都是 \(E_{\max}\)。
当电压升高时，哪层介质先击穿？
两筒间所加的最大电势差是多少？

\paragraph{习题 15.4}
一平板电容器两极板间充满相对介电常量为
\(\varepsilon_r\) 的电介质，电介质两面带电量为
\(Q\)。证明：与金属板相靠的电介质表面所带的
面束缚电荷的电量为
\[
 Q'=\left(1-\frac1{\varepsilon_r}\right)Q.
\]
\textit{这里 \(Q'\) 表示束缚电荷大小，符号与相邻自由面电荷相反。}

\paragraph{习题 15.5}
空气的介电强度为 \(3\,\mathrm{kV/mm}\)，
试求空气中半径分别为
\(1.0\,\mathrm{cm},1.0\,\mathrm{mm},0.1\,\mathrm{mm}\)
的长直导线上单位长度最多各能带多少电荷？

\paragraph{习题 15.6}
人体的某些细胞壁两侧带有等量异号电荷。
设某细胞壁厚 \(5.2\times10^{-9}\,\mathrm m\)，
两表面面电荷密度为
\(\pm0.52\times10^{-3}\,\mathrm{C/m^2}\)
（外表面为正）。
若细胞壁介质相对介电常量为 \(6.0\)，求：
（1）细胞壁内的场强；
（2）细胞壁两表面间的电势差。

\paragraph{习题 15.7}
一块大的均匀电介质平板放在电场强度
\(E_0\) 的均匀电场中，电场方向与板的夹角为
\(\theta\)，如原书图15.20。
已知板的相对介电常量为 \(\varepsilon_r\)，
求板面的面束缚电荷密度。
\textit{图示：} \(\theta\) 为电场与板\emph{平面}之间的夹角，
而非与板法线的夹角。

\paragraph{习题 15.8}
有的计算机键盘的每一个键下面连一小块金属片，
它下面隔一定空气隙是一块小的固定金属片，
两金属片组成小电容器（图15.21）。
按键时，电容改变，电子线路检测后产生信号。
设每个金属片面积 \(50.0\,\mathrm{mm^2}\)，
两金属片间距为 \(0.600\,\mathrm{mm}\)。
如果线路能检出电容变化 \(0.250\,\mathrm{pF}\)，
那么按键需下压多少距离才能给出信号？

\paragraph{习题 15.9}
用两面夹有铝箔的厚为
\(5\times10^{-2}\,\mathrm{mm}\)，
相对介电常量为 \(2.3\) 的聚乙烯膜做一电容器。
如果电容为 \(3.0\,\mu\mathrm F\)，
则膜的面积要多大？

\paragraph{习题 15.10}
空气的击穿场强为 \(3\times10^3\,\mathrm{kV/m}\)。
当一个平行板电容器两极板间是空气
（而电势差为 \(50\,\mathrm{kV}\)）时，
每平方米面积的电容最大是多少？

\paragraph{习题 15.11}
范德格拉夫静电加速器的球形电极半径为
\(18\,\mathrm{cm}\)。
（1）这个球的电容多大？
（2）为了使它的电势升到
\(2.0\times10^5\,\mathrm V\)，
需给它带多少电量？

\paragraph{习题 15.12}
盖革计数管由一根细金属丝和包围它的
同轴导电圆筒组成。
丝直径 \(2.5\times10^{-2}\,\mathrm{mm}\)，
圆筒内直径 \(25\,\mathrm{mm}\)，
管长 \(100\,\mathrm{mm}\)，
管内为真空。
计算盖革计数管的电容
（可用无限长导体圆筒的场强公式计算电场）。

\paragraph{习题 15.13}
图15.22示为调频收音机中的一种可变空气电容器。
奇数极板和偶数极板分别连在一起，其中一组固定，
另一组可转动。
假设极板总数为 \(n\)，每块极板面积为 \(S\)，
相邻两极板距离为 \(d\)，
证明其最大电容为
\[
 C_{\max}=\frac{(n-1)\varepsilon_0S}{d}.
\]

\paragraph{习题 15.14}
一个平行板电容器每个极板面积
\(0.02\,\mathrm{m^2}\)，两板相距 \(0.5\,\mathrm{mm}\)，
放在金属盒中（图15.23）。
电容器两板至盒上下底面的距离各为
\(0.25\,\mathrm{mm}\)。
忽略边缘效应，求此电容器的电容。
若把一个极板与金属盒用导线连接起来，
电容器的电容又是多少？

\paragraph{习题 15.15}
一个电容器由两块长方形金属平板组成
（图15.24），两板长度为 \(a\)、宽度为 \(b\)。
两宽边互相平行，两长边的一端相距 \(d\)，
另一端略微抬起一段距离 \(l\)（\(l\ll d\)）。
板间为真空，求电容。
\textit{图示：}沿板长距离 \(x\in[0,a]\)，
板间距线性变化 \(h(x)=d+(l/a)x\)。

\paragraph{习题 15.16}
为测量电介质材料的相对介电常量，
将一块厚为 \(1.5\,\mathrm{cm}\) 的平板材料
慢慢地插进一个电容器的距离为
\(2.0\,\mathrm{cm}\) 的两平行极板之间。
在插入过程中电容器的电荷保持不变，
插入之后，两板间的电势差减小为原来的
\(60\%\)。
求电介质的相对介电常量。

\paragraph{习题 15.17}
两个同心导体球壳内、外球壳半径分别为
\(R_1,R_2\)，求两者组成的电容器的电容。
把 \(\Delta R=R_2-R_1\ll R_1\) 的极限
与平行板电容器的电容作比较。

\paragraph{习题 15.18}
将一个 \(12\,\mu\mathrm F\) 和两个
\(2\,\mu\mathrm F\) 的电容器连接起来，
组成电容为 \(3\,\mu\mathrm F\) 的电容器组。
如果每个电容器的击穿电压都是
\(200\,\mathrm V\)，
则此电容器组能承受的最大电压是多少？

\paragraph{习题 15.19}
一平行板电容器面积为 \(S\)，
板间距离为 \(d\)，
板间以两层厚度相同、相对介电常量分别为
\(\varepsilon_{r1},\varepsilon_{r2}\)
的电介质充满（图15.25）。
求此电容器的电容。

\paragraph{习题 15.20}
一种利用电容器测量油箱中油量的装置示意如图15.26。
电路测出等效相对介电常量
\(\varepsilon_{r,\mathrm{eff}}\)
（相当于两板间充满板间介质的相对介电常量）。
设电容器两板高度都是 \(a\)，
试导出等效相对介电常量与油面高度
\(h\) 的关系（油相对介电常量为 \(\varepsilon_r\)）。
比较汽油 \(\varepsilon_r=1.95\)
与甲醇 \(\varepsilon_r=33\)：
哪种燃料更宜用这种油量计？

\paragraph{习题 15.21}
一球形电容器的两球间下半部充满
相对介电常量为 \(\varepsilon_r\) 的油，
它的电容较未充油前变化了多少？

\paragraph{习题 15.22}
一个 \(4\,\mu\mathrm F\) 的电容器和一个
\(6\,\mu\mathrm F\) 的电容器串联，
接到 \(200\,\mathrm V\) 的电源上。
充电后断开电源并将两电容器分离。
在下列两种情况下，每个电容器的电压各变为多少？
（1）把每个电容器的正极与另一个电容器的负极相连；
（2）把两电容器的正极与正极相连，负极与负极相连。

\paragraph{习题 15.23}
把一个 \(100\,\mathrm{pF}\) 的电容器充到
\(100\,\mathrm V\)，然后与电源断开，
再把它与另一电容器并联，最后电压为
\(30\,\mathrm V\)。
第二个电容器的电容多大？
并联时损失了多少电能？这电能哪里去了？

\paragraph{习题 15.24}
一个平行板电容器极板面积 \(S\)，
板间距 \(d\)（图15.27）。
（1）充电后保持电量 \(Q\) 不变，
将厚度为 \(b\) 的金属板平行于两极板插入，
电容器储能增加多少？
（2）导体板进入时，外力（非电力）对它做功多少？
是被吸入还是需要推入？
（3）如果充电后使电容器的电压 \(U\) 不变，
则（1）、（2）的结果如何？
\textit{图示：}插入板不接触任何极板，
忽略边缘场，板间有效无场距离变为 \(d-b\)。

\paragraph{习题 15.25}
如图15.28，桌面上固定一半径为
\(7\,\mathrm{cm}\) 的金属圆筒，
其中共轴地吊着一半径为
\(5\,\mathrm{cm}\) 的另一金属圆筒。
今将两筒间加 \(5\,\mathrm{kV}\) 电压，
再移去电源。
求内筒受向下的电力（注意利用功能关系）。

\paragraph{习题 15.26}
一平行板电容器极板长为 \(a\)、
宽为 \(b\)，两板相距为 \(\delta\)
（图15.29）。
对它充电使带电量为 \(Q\)，再把电源断开。
（1）两板间为真空时，电容器储存的电能是多少？
（2）板间插入宽为 \(b\)、厚为 \(\delta\)、
相对介电常量为 \(\varepsilon_r\) 的均匀介质板，
插入距离为 \(x\)，储存电能是多少？
（3）介质板插入距离 \(x\) 时，
受到电力的大小和方向如何？
（4）当 \(x\) 从0增加至 \(a\) 时，
能量如何转化（设无摩擦）？
（5）若插板过程中两极板保持与电压
恒定为 \(U\) 的电源相连，
重新解答（2）—（4），注意电源对系统做功。

\paragraph{习题 15.27}
证明：球形电容器带电后，
其电场能量一半储存在内半径为
\(R_1\)、外半径为
\(2R_1R_2/(R_1+R_2)\) 的球壳内，
其中 \(R_1,R_2\) 分别为电容器的内外球半径。
对于一个孤立导体球，带电后电场能量的一半
储存在多大的球壳内？

\paragraph{习题 15.28}
一个平行板电容器极板面积 \(S\)，
板间距 \(y_0\)，下极板在 \(y=0\)，
上极板在 \(y=y_0\)。
充满两板间的介质相对介电常量沿 \(y\) 变化，
关系为
\[
 \varepsilon_r(y)=1+\frac{3y}{y_0}.
\]
（1）求电容；
（2）上极板带 \(+Q\) 时介质上、下表面的面束缚电荷密度；
（3）用高斯定律求介质体内束缚电荷密度；
（4）证明体束缚电荷与上下表面束缚电荷的总和为零。

\paragraph{习题 15.29}
一个中空铜球浮在相对介电常量为
\(3.0\) 的大油缸中，一半没入油内。
若铜球总电量为
\(2.0\times10^{-6}\,\mathrm C\)，
它的上半部与下半部各带多少电量？

\paragraph{习题 15.30}
在具有杂质离子的半导体中，
电子围绕这些离子作轨道运动。
若轨道尺度大于半导体原子间距，
可视为电子在介电常量近似均匀的介质空间中运动。
（1）按习题13.34描述的玻尔理论，
计算一个电子的轨道能；
（2）若半导体相对介电常量为
\(\varepsilon_r=15.8\)，估计电子在最小玻尔轨道中的
运动能量，与真空中电子绕离子运动的最小轨道能相比如何？
